% GENERATED FILE. Edit canonical lessons, then run npm run generate:books.
% canonical-source-hash: fnv1a64:ccde1758516546b5
% canonical-lessons: PA-C250-handi, PA-C250-karaha, PA-C250-garva, PA-C250-gagar, PA-C250-pittal

\chapter{Pots and Pans}
\label{ch:pa-a2-pots-and-pans}
\hlchaptermodality{250}

\begin{chapteropening}
\textbf{By the end of this chapter:} \emph{I can name cooking pots and what they are made of.}

Complete the last lesson of chapter 250: I can name cooking pots and what they are made of.
\end{chapteropening}

\section[\texorpdfstring{\pa{ਹਾਂਡੀ} (hāṇḍī)}{ਹਾਂਡੀ (hāṇḍī)}]{\pa{ਹਾਂਡੀ} (hāṇḍī) — a clay cooking pot}

\label{lesson:PA-C250-handi}



\begin{quote}
Before the new one: say the Punjabi for charcoal, then the Punjabi for an ember.
\end{quote}

\subsection*{You'll want to know: \pa{ਹਾਂਡੀ}}
\textbf{\pa{ਹਾਂਡੀ}} — \emph{hāṇḍī} — \textquotedblleft{}a clay cooking pot\textquotedblright{}.

\begin{quote}
\pa{ਹਾਂਡੀ} \pa{ਚੁੱਲ੍ਹੇ} \pa{ਉੱਤੇ} \pa{ਹੈ।}
\end{quote}

\begin{quote}
\emph{hāṇḍī chullhe utte hai.}
\end{quote}

\begin{quote}
The clay pot is on the stove.
\end{quote}

A feminine noun.

\begin{cousinweb}
\textbf{\pa{ਹਾਂਡੀ}} (\emph{hāṇḍī}) is inherited from a reconstructed \emph{*hāṇḍa}, \textquotedblleft{}pot\textquotedblright{}; CDIAL lists Punjabi \emph{hā̃ḍī}, \textquotedblleft{}cooking pot\textquotedblright{} (CDIAL 14050).
\end{cousinweb}

\begin{grammarlens}[title={What you've built}]
One more word for this chapter.
\end{grammarlens}

\subsection*{Your turn}
\begin{itemize}
\raggedright
  \item \emph{Say it:} \emph{hāṇḍī}
  \item \emph{Say it:} \emph{hāṇḍī}, once more
  \item \emph{Say it:} \emph{hāṇḍī chullhe utte hai.} — The clay pot is on the stove.
  \item \emph{Recall:} say the Punjabi for charcoal, then the Punjabi for an ember, then say \emph{hāṇḍī} again
\end{itemize}

\begin{tcolorbox}[breakable,colback=teal!4,colframe=teal!35!black,title={Before you move on}]
What does \pa{ਹਾਂਡੀ} mean? (\textquotedblleft{}A clay cooking pot\textquotedblright{}.) Say it once more.
\end{tcolorbox}
\section[\texorpdfstring{\pa{ਕੜਾਹਾ} (kaṛāhā)}{ਕੜਾਹਾ (kaṛāhā)}]{\pa{ਕੜਾਹਾ} (kaṛāhā) — a large pan}

\label{lesson:PA-C250-karaha}



\begin{quote}
Before the new one: say the Punjabi for an ember, then the Punjabi for a clay cooking pot.
\end{quote}

\subsection*{You'll want to know: \pa{ਕੜਾਹਾ}}
\textbf{\pa{ਕੜਾਹਾ}} — \emph{kaṛāhā} — \textquotedblleft{}a large pan\textquotedblright{}.

\begin{quote}
\pa{ਕੜਾਹੇ} \pa{ਵਿੱਚ} \pa{ਬਹੁਤ} \pa{ਹਲਵਾ} \pa{ਹੈ।}
\end{quote}

\begin{quote}
\emph{kaṛāhe vich bahut halvā hai.}
\end{quote}

\begin{quote}
There is a lot of halwa in the big pan.
\end{quote}

A masculine noun.

\begin{cousinweb}
\textbf{\pa{ਕੜਾਹਾ}} (\emph{kaṛāhā}) is inherited from Sanskrit \emph{kaṭāha}, \textquotedblleft{}boiler, saucepan with handles\textquotedblright{}; CDIAL lists Punjabi \emph{kaṛāhā}, \textquotedblleft{}shallow pan for boiling sugarcane juice\textquotedblright{} (CDIAL 2638).
\end{cousinweb}

\begin{grammarlens}[title={What you've built}]
One more word for this chapter.
\end{grammarlens}

\subsection*{Your turn}
\begin{itemize}
\raggedright
  \item \emph{Say it:} \emph{kaṛāhā}
  \item \emph{Say it:} \emph{kaṛāhā}, once more
  \item \emph{Say it:} \emph{kaṛāhe vich bahut halvā hai.} — There is a lot of halwa in the big pan.
  \item \emph{Recall:} say the Punjabi for an ember, then the Punjabi for a clay cooking pot, then say \emph{kaṛāhā} again
  \item \emph{Write it:} copy \textbf{\pa{ਕੜਾਹਾ}} once, checking the whole word
\end{itemize}

\begin{tcolorbox}[breakable,colback=teal!4,colframe=teal!35!black,title={Before you move on}]
What does \pa{ਕੜਾਹਾ} mean? (\textquotedblleft{}A large pan\textquotedblright{}.) Say it once more.
\end{tcolorbox}
\section[\texorpdfstring{\pa{ਗੜਵਾ} (gaṛvā)}{ਗੜਵਾ (gaṛvā)}]{\pa{ਗੜਵਾ} (gaṛvā) — a metal jug}

\label{lesson:PA-C250-garva}



\begin{quote}
Before the new one: say the Punjabi for a clay cooking pot, then the Punjabi for a large pan.
\end{quote}

\subsection*{You'll want to know: \pa{ਗੜਵਾ}}
\textbf{\pa{ਗੜਵਾ}} — \emph{gaṛvā} — \textquotedblleft{}a metal jug\textquotedblright{}.

\begin{quote}
\pa{ਗੜਵੇ} \pa{ਵਿੱਚ} \pa{ਲੱਸੀ} \pa{ਹੈ।}
\end{quote}

\begin{quote}
\emph{gaṛve vich lassī hai.}
\end{quote}

\begin{quote}
There is lassi in the jug.
\end{quote}

A masculine noun.

\begin{cousinweb}
\textbf{\pa{ਗੜਵਾ}} (\emph{gaṛvā}) is inherited from Sanskrit \emph{gaḍḍuka}, \textquotedblleft{}waterpot\textquotedblright{}; CDIAL lists Punjabi \emph{gaṛvā}, \textquotedblleft{}brass jug\textquotedblright{} (CDIAL 3984). CDIAL records the view that the Sanskrit word may be from a Munda language.
\end{cousinweb}

\begin{grammarlens}[title={What you've built}]
One more word for this chapter.
\end{grammarlens}

\subsection*{Your turn}
\begin{itemize}
\raggedright
  \item \emph{Say it:} \emph{gaṛvā}
  \item \emph{Say it:} \emph{gaṛvā}, once more
  \item \emph{Say it:} \emph{gaṛve vich lassī hai.} — There is lassi in the jug.
  \item \emph{Recall:} say the Punjabi for a clay cooking pot, then the Punjabi for a large pan, then say \emph{gaṛvā} again
  \item \emph{Write it:} copy \textbf{\pa{ਗੜਵਾ}} once, checking the whole word
\end{itemize}

\begin{tcolorbox}[breakable,colback=teal!4,colframe=teal!35!black,title={Before you move on}]
What does \pa{ਗੜਵਾ} mean? (\textquotedblleft{}A metal jug\textquotedblright{}.) Say it once more.
\end{tcolorbox}
\section[\texorpdfstring{\pa{ਗਾਗਰ} (gāgar)}{ਗਾਗਰ (gāgar)}]{\pa{ਗਾਗਰ} (gāgar) — a water pitcher}

\label{lesson:PA-C250-gagar}



\begin{quote}
Before the new one: say the Punjabi for a large pan, then the Punjabi for a metal jug.
\end{quote}

\subsection*{You'll want to know: \pa{ਗਾਗਰ}}
\textbf{\pa{ਗਾਗਰ}} — \emph{gāgar} — \textquotedblleft{}a water pitcher\textquotedblright{}.

\begin{quote}
\pa{ਮੈਂ} \pa{ਗਾਗਰ} \pa{ਪਾਣੀ} \pa{ਨਾਲ} \pa{ਭਰਦੀ} \pa{ਹਾਂ।}
\end{quote}

\begin{quote}
\emph{maiṁ gāgar pāṇī nāl bhardī hāṁ.}
\end{quote}

\begin{quote}
I fill the pitcher with water.
\end{quote}

A feminine noun.

\begin{cousinweb}
\textbf{\pa{ਗਾਗਰ}} (\emph{gāgar}) is inherited from Sanskrit \emph{gargara}, \textquotedblleft{}churn, waterpot\textquotedblright{}; CDIAL lists Punjabi \emph{gāgar}, \textquotedblleft{}waterpot\textquotedblright{} (CDIAL 4043).
\end{cousinweb}

\begin{grammarlens}[title={What you've built}]
One more word for this chapter.
\end{grammarlens}

\subsection*{Your turn}
\begin{itemize}
\raggedright
  \item \emph{Say it:} \emph{gāgar}
  \item \emph{Say it:} \emph{gāgar}, once more
  \item \emph{Say it:} \emph{maiṁ gāgar pāṇī nāl bhardī hāṁ.} — I fill the pitcher with water.
  \item \emph{Recall:} say the Punjabi for a large pan, then the Punjabi for a metal jug, then say \emph{gāgar} again
  \item \emph{Write it:} copy \textbf{\pa{ਗਾਗਰ}} once, checking the whole word
\end{itemize}

\begin{tcolorbox}[breakable,colback=teal!4,colframe=teal!35!black,title={Before you move on}]
What does \pa{ਗਾਗਰ} mean? (\textquotedblleft{}A water pitcher\textquotedblright{}.) Say it once more.
\end{tcolorbox}
\section[\texorpdfstring{\pa{ਪਿੱਤਲ} (pittal)}{ਪਿੱਤਲ (pittal)}]{\pa{ਪਿੱਤਲ} (pittal) — brass}

\label{lesson:PA-C250-pittal}



\begin{quote}
Before the new one: say the Punjabi for a metal jug, then the Punjabi for a water pitcher.
\end{quote}

\subsection*{You'll want to know: \pa{ਪਿੱਤਲ}}
\textbf{\pa{ਪਿੱਤਲ}} — \emph{pittal} — \textquotedblleft{}brass\textquotedblright{}.

\begin{quote}
\pa{ਇਹ} \pa{ਗੜਵਾ} \pa{ਪਿੱਤਲ} \pa{ਦਾ} \pa{ਹੈ।}
\end{quote}

\begin{quote}
\emph{eh gaṛvā pittal dā hai.}
\end{quote}

\begin{quote}
This jug is made of brass.
\end{quote}

A masculine noun.

\begin{cousinweb}
\textbf{\pa{ਪਿੱਤਲ}} (\emph{pittal}) is inherited from Sanskrit \emph{pittala}, \textquotedblleft{}brass\textquotedblright{}; CDIAL lists Punjabi \emph{pittal}, \textquotedblleft{}brass\textquotedblright{} (CDIAL 8184).
\end{cousinweb}

\begin{grammarlens}[title={What you've built}]
One more word for this chapter.
\end{grammarlens}

\subsection*{Your turn}
\begin{itemize}
\raggedright
  \item \emph{Say it:} \emph{pittal}
  \item \emph{Say it:} \emph{pittal}, once more
  \item \emph{Say it:} \emph{eh gaṛvā pittal dā hai.} — This jug is made of brass.
  \item \emph{Recall:} say the Punjabi for a metal jug, then the Punjabi for a water pitcher, then say \emph{pittal} again
  \item \emph{Write it:} copy \textbf{\pa{ਪਿੱਤਲ}} once, checking the whole word
\end{itemize}

\begin{tcolorbox}[breakable,colback=teal!4,colframe=teal!35!black,title={Before you move on}]
What does \pa{ਪਿੱਤਲ} mean? (\textquotedblleft{}Brass\textquotedblright{}.) Say it once more.
\end{tcolorbox}
